\section{Methodology}
\label{sec:methodology}

This section presents our methodology for Python-anchored cross-language code
clone detection (XCCD) using execution-augmented search, inspired by
rStar-Math. The central idea is to transform clone detection from a direct
classification task into a search-and-verification problem: given two
programs, one written in Python and the other in some other programming
language, a policy model is required to reason about their input/output
behavior, decide whether direct execution of both programs is warranted, run
both programs when it is, and produce a stable, extractable conclusion about
whether they are semantic clones. Multiple candidate reasoning trajectories
are generated per pair using Monte Carlo Tree Search (MCTS), and a dedicated
evaluator aggregates evidence across the resulting search tree into a single
prediction. No model training, fine-tuning, or reinforcement learning is
performed in the work described here; the framework is designed so that the
trajectories it produces can later be used as training data, but that
extension is left to future work.

Our methodology has four main contributions. First, we design a structured
prompt that forces the policy model to produce a stable and verifiable
clone-detection output, expressed as an explicit reasoning-and-code trace
rather than a bare label. Second, we let the model execute \emph{both}
programs directly -- the Python side as an interpreted subprocess and the
non-Python side compiled or interpreted as appropriate for its language --
rather than asking the model to simulate or predict the second program's
behavior. Third, we use MCTS as a search mechanism that generates and scores
multiple candidate clone-analysis trajectories per pair, exploring alternative
reasoning and testing strategies rather than committing to a single greedy
generation. Fourth, we design a dedicated evaluator that aggregates the
(possibly several, possibly disagreeing) leaf-level verdicts produced by a
search tree into a single, calibrated prediction, and that treats a tree's own
inability to settle a comparison as a further evidential signal in its own
right.

\subsection{Scope and Problem Formulation}
\label{subsec:scope_problem}

Let $c_p$ denote a Python source-code snippet and let $c_x$ denote a
source-code snippet written in a non-Python programming language
$\ell_x \in \mathcal{L}$, where $\mathcal{L}$ is the set of target languages
under study (e.g., Java, Rust). In this study, we focus on Python-anchored
XCCD, where every evaluated pair contains one Python program:
\[
    x = (c_p, c_x),
    \qquad
    \ell(c_p)=\text{Python},
    \qquad
    \ell(c_x)=\ell_x \neq \text{Python}.
\]
The goal is to predict a binary clone label
\[
    y \in \{\text{clone}, \text{non-clone}\}.
\]

This Python-anchored setting is a deliberate scope restriction that lets us
evaluate the core hypothesis of this work -- whether letting a language model
generate and directly execute behavioral tests improves XCCD -- without first
solving fully language-agnostic XCCD (e.g., Java--Ruby or C++--Rust). Because
one side of every pair is Python, the policy model can always ground its
analysis in a directly runnable, interpreted reference implementation, while
the non-Python side is executed through the appropriate toolchain for
$\ell_x$ (Section~\ref{subsec:execution_model}).

Two programs are considered semantic clones if they accept compatible input
formats and produce equivalent outputs for valid inputs in the shared input
domain. We define the input signature of a program $c$ as
\[
    \sigma(c) = (\tau, d, \kappa),
\]
where $\tau$ represents the input data types, $d$ represents the input
dimensionality or structure, and $\kappa$ represents additional constraints
such as input size, order, and formatting. The policy model's own reasoning
step (Section~\ref{subsec:prompt_design}) begins by estimating and comparing
$\sigma(c_p)$ and $\sigma(c_x)$. If it judges
$\sigma(c_p) \neq \sigma(c_x)$ -- i.e., the two programs differ at the
signature level, or are plainly unrelated -- it is instructed to skip
behavioral testing and emit \texttt{print("not code clones")} directly,
concluding non-clone without executing either program. Otherwise, it proceeds
to generate and run behavioral tests before concluding.

\paragraph{Remark on signature-based pruning.}
Treating an apparent signature mismatch as sufficient evidence of non-clone is
an operational simplification, not a guarantee: two genuine clones may
legitimately differ in surface input conventions (for instance, reading three
integers on separate lines versus a single comma-separated line), so an overly
literal signature comparison risks false negatives. Unlike a hand-engineered
canonicalization pass, this comparison is made by the policy model itself as
part of its own reasoning trace, guided by the prompt's explicit instruction
that a merely \emph{suspected} behavioral difference -- as opposed to a
genuine difference in input/output type, arity, or purpose -- is not
sufficient grounds to skip execution and must instead be settled by actually
running both programs. This keeps signature-based pruning available as a
cheap early exit for clearly unrelated pairs while pushing every borderline
case toward direct execution.

\subsection{Prompt Design}
\label{subsec:prompt_design}

The first contribution of this work is a structured prompt that requires the
policy model to produce its analysis as an explicit, machine-parsable
trajectory rather than a bare label. The model's response is required to
consist of four tagged segments -- \texttt{<analysis>}, \texttt{<code>},
\texttt{<output>}, and \texttt{<answer>} -- with a final decision expressed
as \verb|\boxed{clone}| or \verb|\boxed{non-clone}| inside the answer segment.
The prompt used in this work is shown in
Fig.~\ref{fig:xccd-prompt}; it is language-agnostic modulo the specific
compiler or interpreter invoked for $\ell_x$.

\begin{figure}[t]
\centering
\begin{tcolorbox}[width=0.95\linewidth,
title=\textbf{Prompt for Python-Anchored XCCD}, halign=left]
\ttfamily\small
Analyze the following two code snippets and determine whether they are
semantic code clones. Focus on semantic behavior, not surface-level syntax.

Your solution must include four parts: <analysis>, <code>, <output>, and
<answer>.

In <analysis>, summarize the behavior of each code snippet and compare their
input/output behavior. Consider input type, output type, main logic, edge
cases, and whether both snippets produce equivalent outputs for equivalent
valid inputs.

If the two snippets could plausibly be semantic clones, do not predict what
Code 2 would print. Both snippets are complete programs that read standard
input and write standard output, so run BOTH of them on the same input and
compare the output they actually produce. A decision backed by two observed
outputs is always worth more than a decision backed by reasoning about one of
them.

Code 2 is written in \{language\}. Compile and run it directly, the same way
you run Code 1: the harness saves Code 2's exact source and compiles it
before your step runs, appending the compile status to the output you
receive; invoke the compiled program the same way you invoke Code 1. If
compilation fails, the error text tells you why -- only then fall back to
reasoning about Code 2's expected behavior, and say that the comparison is
based on expected behavior rather than direct execution.

Skip testing ONLY when the two snippets differ at the signature level --
different input types, different output types, or plainly unrelated
functionality. A suspicion that they behave differently is not a reason to
skip: if the difference you suspect is behavioural (an off-by-one, a
rounding or division rule, an edge case, a different tie-break), that is
precisely what running both sides settles, so run them and let the two
outputs decide. In the signature-level case only, write in the <code>
section: print("not code clones")

Keep every <code> step short: a few lines, and at most one or two test
inputs per step. If you want more test cases, put them in additional steps
instead of one long block.

The final decision must be either clone or non-clone. Put the final decision
inside \textbackslash boxed\{\}.

Question: \{input\}
\end{tcolorbox}
\caption{Structured prompt used for Python-anchored XCCD. \{language\} and
\{input\} are filled in per question; the model is additionally given a
small number of worked examples in the same four-segment format, one of
which is a signature-mismatch case resolved without execution and one of
which is a behavioral case resolved by running both programs.}
\Description{}
\label{fig:xccd-prompt}
\end{figure}

A generated trajectory is required to end each reasoning or code segment with
an explicit step delimiter, so that the search procedure
(Section~\ref{subsec:mcts_search}) can cleanly split a response into discrete
steps. The prompt also discourages relying on variable names, comments, or
superficial syntactic similarity, and instructs the model to keep each
\texttt{<code>} step short -- long code blocks that overrun the per-step
generation budget are truncated and discarded, so the model is guided to
spread multiple test cases across several short steps rather than one long
one.

\subsection{Execution Model}
\label{subsec:execution_model}

Unlike a design that only executes the Python side and asks the model to
simulate or predict the non-Python program's output, both programs are
executed directly whenever the model chooses to test. The Python snippet
$c_p$ is saved to a scratch file and invoked as a subprocess of the Python
interpreter. The non-Python snippet $c_x$ is saved to a scratch file named
appropriately for $\ell_x$ (e.g., by its declared class name for Java) and,
for compiled languages, compiled with the corresponding compiler (e.g.,
\texttt{javac} for Java, \texttt{rustc} for Rust) before being invoked; for
interpreted languages it is invoked directly. Compilation is triggered
implicitly the first time the model's generated code names the compiler or
the resulting binary, and the resulting compile status -- success, or the
compiler's error output on failure -- is appended to what the model observes
for that step, so a failed compilation is visible evidence the model can
react to rather than a silent gap. Programs that depend on external libraries
unavailable in the execution environment are reported as such rather than
silently failing, and the model is instructed to fall back to reasoning about
expected behavior only in that case.

Given a test input $t$, executing $c_p$ and $c_x$ yields raw standard-output
text
\[
    o_p = R_{\text{Python}}(c_p,t), \qquad o_x = R_{\ell_x}(c_x,t).
\]
Equivalence between $o_p$ and $o_x$ is judged by the policy model itself,
typically by comparing the two observed outputs directly (optionally after
stripping incidental whitespace) inside the same \texttt{<code>} step that
produced them, rather than by a separate, formally specified cross-language
equivalence procedure external to the model. This keeps the comparison
flexible across arbitrary output shapes but means the model's own judgment of
what counts as "the same output" is itself part of what is being evaluated;
Section~\ref{subsec:limitations} discusses failure modes this can produce,
such as numeric-formatting or integer-overflow discrepancies between
languages that a literal string comparison treats as a genuine mismatch, or
conversely a coincidental match on the specific inputs tested.

All execution is subject to a bounded number of consecutive failures per
trajectory (Section~\ref{subsec:mcts_search}) and runs in an isolated working
directory per search job, so that concurrently generated scratch files and
compiled artifacts from different questions or rollouts cannot interfere with
one another. Environment state inherited from the model-serving process that
would otherwise cause spurious, environment-induced compilation failures for
certain languages is deliberately excluded from the compiler subprocess's
environment.

\subsection{Search Tree Construction}
\label{subsec:search_space}

Following the rStar-Math style of step-by-step reasoning, a clone-analysis
response is decomposed into a sequence of reasoning and code steps. A
reasoning trajectory is
\[
    T_r = (s_1, s_2, \ldots, s_m),
\]
where each $s_i$ is one step of model-generated text -- a chunk of
\texttt{<analysis>} reasoning, a \texttt{<code>} block together with its
appended \texttt{<output>} observation, or the terminal \texttt{<answer>}
segment -- ending at a step delimiter. Unlike a fixed, hand-specified
decomposition, $m$ is not fixed in advance: the prompt allows (and, for
multi-test-case reasoning, encourages) the number of intermediate steps to
vary per trajectory, subject only to a maximum search depth.

Each step is a node in a search tree. The root node is the original pair
$(c_p,c_x)$, and child nodes are generated by extending the trajectory up to
that node with a newly sampled continuation. A node is terminal if it
contains a parsed final answer, or if the trajectory has failed -- through
too many consecutive code-execution errors, exceeding the maximum depth, or
producing text that cannot be parsed into a valid step -- in which case it is
assigned a fixed failure marker rather than a clone/non-clone label.

\subsection{MCTS-Based Search}
\label{subsec:mcts_search}

MCTS is used purely as an inference-time search and verification mechanism:
for each pair $(c_p,c_x)$, the policy model $\pi_\theta$ generates multiple
candidate trajectories, which are then aggregated into a single prediction by
the evaluator (Section~\ref{subsec:evaluator}). No model parameters are
updated during this process.

At a non-terminal node $s$ with at least one non-terminal child, the next
child to expand from is chosen by an upper-confidence selection rule:
\[
    \operatorname{UCT}(s)
    =
    Q(s)
    +
    c_{\text{puct}}
    \frac{\sqrt{N(\operatorname{parent}(s))}}{1 + N(s)},
\]
where $Q(s)$ is the running average value backed up to $s$, $N(s)$ is its
visit count, and $c_{\text{puct}}$ controls the exploration-exploitation
trade-off. $Q(s)$ is treated as zero for an as-yet-unvisited node, but the
exploration term is not: it is well-defined and strictly positive whenever
the parent itself has been visited, so an unvisited child is always credited
for being unexplored rather than scored identically to a sibling that has
already accumulated a poor outcome. Expansion itself proceeds by sampling
several candidate continuations from the policy model at temperature $>0$
for the selected node, each of which becomes a new child; the number of
continuations sampled per expansion is a configuration parameter of the
search (Section~\ref{subsec:experimental_setting}).

\paragraph{Technical contribution: correcting a degenerate exploration term.}
An earlier implementation of the selection rule above set the exploration
term to zero not only for $Q(s)$ but for the entire $u$-term whenever
$N(s)=0$, collapsing $\operatorname{UCT}(s)$ to exactly $0$ for every
unvisited child regardless of $c_{\text{puct}}$ or how many times the parent
had already been visited. Because every child of a freshly expanded node
starts at $N(s)=0$, this made the first-listed child of an expansion -- an
artifact of implementation order, not a property of the trajectory itself --
strictly preferred over its siblings the moment it received any positive
backed-up value: a positive $Q(s)$ then permanently outranked the
$\operatorname{UCT}(s)=0$ score every unvisited sibling was stuck at, with no
mechanism for the score to ever recover. In practice this let the search
spend its entire remaining rollout budget at a node re-descending a single
child's subtree, while the other $n_{\text{generate\_sample}}-1$ children
produced at that expansion -- each already the product of a real
policy-model sample and, whenever its step involved code, a real execution
or compilation -- were never selected again. The failure is bounded at the
default branching factor of two candidates per expansion
(Section~\ref{subsec:experimental_setting}), where at most one sibling can be
starved, but compounds with wider expansions: at
$n_{\text{generate\_sample}}=4$, for instance, we observed the search
degenerate into repeatedly extending one branch step after step until the
maximum depth was reached, rather than exploring the alternative
continuations it had already paid to generate -- behavior that is
functionally indistinguishable from an infinite loop from the outside, since
wall-clock cost scales with the width of every expansion while the effective
search width collapses to one.

We apply two corrections, verified in isolation to compose safely (an
unvisited child is always selected ahead of a visited one, and the first full
pass over a node's children proceeds in a fixed, non-repeating order). First,
the $u$-term is rewritten as $c_{\text{puct}}\sqrt{N(\operatorname{parent}(s))}
/ (1+N(s))$, the standard AlphaZero-style form: it requires no special case
for $N(s)=0$, evaluates to a strictly positive bonus for an unvisited child
as soon as its parent has been visited even once, and still decays as the
child itself accumulates visits, so it does not prevent the search from
eventually converging on exploiting a genuinely strong child. Second, as a
supplementary, threshold-style guard layered on top of this corrected term,
child selection at a node restricts competition to that node's not-yet-visited
children whenever at least one exists, only falling back to comparing
$\operatorname{UCT}(s)$ across all non-terminal children once every child has
received at least one visit. This guarantees every candidate continuation
produced at an expansion is actually explored once before any one of them is
exploited a second time, independent of how the corrected exploration bonus
happens to be tuned via $c_{\text{puct}}$. The guard is not free: clearing one
level of the tree under it requires at least $n_{\text{generate\_sample}}$
rollouts, so for configurations where the rollout budget is smaller than the
branching factor, the guard itself becomes the binding constraint on how deep
the search can reach -- we flag this interaction explicitly rather than
treating the guard as strictly dominant over the corrected exploration term
on its own.

Backup along a trajectory is driven by execution outcomes rather than by a
learned process-reward model in the setting evaluated in this work: when a
generated code step raises a genuine execution error (as opposed to, e.g., a
self-written assertion the model used to test its own equivalence claim,
which is scored separately for internal consistency once the branch's
eventual conclusion is known), a negative value is backed up from that node
to the root, discouraging the search from re-exploring branches that produce
broken code. Ground-truth-informed backup -- in which a trajectory's value is
determined by whether its final conclusion matches the true label -- is
supported by the underlying search implementation for future data-generation
use, but is never invoked in the evaluation setting studied here: the search
does not observe the label at any point while producing its predictions, and
no learned value or reward model is queried during node selection in this
work. A configurable number of independent rollouts is run from the root per
question, and a configurable maximum depth bounds how many steps any single
trajectory may take before it is forced to terminate.

\subsection{The Evaluator: Multi-Trajectory Label Aggregation}
\label{subsec:evaluator}

Because MCTS run under the search procedure above can, and typically does,
produce more than one terminal trajectory per question -- and because those
trajectories can disagree -- our third contribution, and a necessary
component of the overall pipeline, is a dedicated evaluator that turns a
finished search tree into a single prediction. This is a deliberate departure
from selecting a single highest-scoring trajectory: rather than trusting one
path chosen by a value estimate, the evaluator treats every terminal leaf of
the tree as a vote and combines the votes explicitly.

\paragraph{Label extraction.} Every terminal node whose \texttt{<answer>}
segment contains a parsed \verb|\boxed{}| decision contributes one vote. A
leaf's raw decision text is normalized to \texttt{clone} or \texttt{non-clone}
before voting, so that paraphrases of the intended label (for example,
"semantic clone" or "not a clone") are still counted rather than discarded.
Leaves whose trajectory ended in one of the fixed failure markers
(Section~\ref{subsec:search_space}) do not vote.

\paragraph{Vote aggregation.} Given the set of votes $V = \{v_1, \ldots,
v_k\}$ collected from a tree's leaves, the predicted label is
\[
    \hat{y} =
    \begin{cases}
        \text{clone}, & \text{if } \exists\, v_i = \text{clone} \text{ (clone-on-disagreement)}, \\[4pt]
        \displaystyle\arg\max_{c \in \{\text{clone},\text{non-clone}\}}
        \bigl|\{v_i \in V : v_i = c\}\bigr|,
        & \text{under plain majority voting},
    \end{cases}
\]
with ties under plain majority broken in favor of the earliest-generated
leaf. Which of the two aggregation rules is used is itself a modeling
decision rather than a fixed default: clone-on-disagreement is only
appropriate when the model's precision on the clone side substantially
exceeds its recall, since conceding any disagreement to "clone" trades away
a small amount of precision to recover false negatives, and is a net loss
whenever that precision advantage does not hold. We select between the two
rules per language pair based on this asymmetry, measured on
held-out data.

\paragraph{Fallback for trees with no valid vote.} If a tree produces no
leaf with a parsable clone/non-clone decision at all -- $V = \emptyset$ -- the
question would otherwise have no prediction. Rather than leaving it
unanswered, the evaluator falls back to a signal drawn from \emph{why} the
tree failed to settle the comparison: if any node in the tree shows a
signature of having been unable to execute anything and complete the
comparison it attempted -- for instance, code that never became a runnable
program, a language side the execution environment does not support, or a
self-raised assertion that was never resolved into a final decision -- the
fallback predicts \texttt{clone}. This is deliberately restricted to trees
with no other evidence: it is a fallback for an otherwise-unanswered
question, not an override of a tree that already reached a real verdict. The
justification is empirical rather than a priori: across held-out evaluation
data, trees exhibiting this signature skew overwhelmingly toward genuinely
being clones, consistent with the intuition that a pair whose equivalence is
hard enough that no branch of the search could pin it down by actually
running code is disproportionately likely to be a case where the two
programs really are equivalent. The specific set of failure signatures that
triggers this fallback is selected by validating candidate trigger subsets on
held-out data rather than by reading whole-run accuracy, to avoid overfitting
the rule to the evaluation set it is meant to improve.

\subsection{Final Decision Rule}
\label{subsec:final_decision}

Putting the search and evaluator together, the label reported for a pair
$(c_p,c_x)$ is
\[
    \hat{y}(c_p,c_x) = \operatorname{Evaluator}\bigl(\mathcal{B}(c_p,c_x)\bigr),
\]
where $\mathcal{B}(c_p,c_x)$ is the finished search tree produced by MCTS for
that pair and $\operatorname{Evaluator}(\cdot)$ applies label extraction,
vote aggregation, and the no-vote fallback as defined in
Section~\ref{subsec:evaluator}. A question for which the evaluator returns no
prediction at all -- which, given the fallback above, occurs only for trees
with no vote and no qualifying failure signature -- is counted separately as
a non-response rather than assigned an arbitrary label.

\subsection{Experimental Setting}
\label{subsec:experimental_setting}

\paragraph{Datasets.} We evaluate on Python-anchored cross-language clone
pairs derived from the CLCCD benchmark, which itself draws on Project
CodeNet and XLCoST. The Python--Java pairs are pooled directly from both
CodeNet- and XLCoST-derived CLCCD records that pair Java with Python,
yielding 1{,}800 pairs balanced 900 clone / 900 non-clone. For target
languages not natively paired with Python in CLCCD (e.g., Rust), pairs are
constructed by joining CodeNet-derived Python and target-language solutions
that solve the same underlying problem (yielding a clone pair) or different
problems (yielding a non-clone pair); this bridged construction is limited by
the number of CodeNet problems with both a Python and a target-language
solution, and non-clone pairs are oversampled from that fixed pool of
problems to reach a usable set size, giving an imbalanced but fixed
clone:non-clone ratio (300 clone / 900 non-clone for the Rust pairs used in
this work).

\paragraph{Model.} We fix the policy model to Qwen3-4B, used without any
additional fine-tuning, rather than evaluating the framework model-agnostically
across a pool of instruction-tuned models of varying size and code
specialization. This choice is the result of a preliminary comparison against
Qwen2.5-Coder-3B-Instruct -- a code-specialized model at a comparable parameter
count -- run under the same search configuration described below and reported
in Table~\ref{tab:pure-vs-mcts}. Qwen3-4B is the model for which MCTS-driven
search reliably improves on the same model's own pure (single-pass,
non-search) inference baseline across both target languages studied here
(Python--Java $F_1$: 0.8767 $\rightarrow$ 0.9454; Python--Rust $F_1$: 0.4806
$\rightarrow$ 0.8739), including a large gain on the harder Python--Rust pair,
whereas Qwen2.5-Coder-3B-Instruct is already close to ceiling on Python--Java
without search (pure-inference $F_1$ 0.9457) and shows a net $F_1$ regression
once the same MCTS configuration is applied there ($F_1$ 0.9226), driven by a
recall loss that its precision gain does not offset. Because a central claim
of this work is that search and direct execution improve XCCD over
single-pass inference, we adopt the model for which that improvement actually
materializes across both languages, rather than the model with the strongest
code-specific pretraining or the highest pure-inference score in isolation.

\paragraph{Search and decoding configuration.} Unless otherwise noted, the
search is run with the following configuration for Qwen3-4B across every
language pair. Sampling temperature 0.7, nucleus sampling $p=0.8$, top-$k=20$,
repetition penalty 1.1, a minimum of 8 generated tokens and a maximum of 512
generated tokens per step. The search explores up to a maximum depth of 16
steps, samples 2 candidate continuations per expansion, and runs 2 independent
rollouts per question. The exploration constant is $c_{\text{puct}}=2$. A
trajectory is terminated as failed after 2 consecutive code-execution errors.
Each prompt includes 2 worked few-shot examples in the same four-segment
format as Fig.~\ref{fig:xccd-prompt}. The search never observes the
ground-truth label at any point while producing its predictions.

\paragraph{Metrics.} We report precision, recall, and $F_1$ over judged
pairs,
\[
    \text{Precision} = \frac{TP}{TP+FP}, \qquad
    \text{Recall} = \frac{TP}{TP+FN}, \qquad
    F_1 = \frac{2\cdot\text{Precision}\cdot\text{Recall}}{\text{Precision}+\text{Recall}},
\]
alongside the \emph{response rate}: the fraction of pairs for which the
evaluator (Section~\ref{subsec:evaluator}) returns a prediction at all,
distinguishing genuine non-response from a wrong-but-confident answer.

\subsection{Limitations}
\label{subsec:limitations}

The Python-anchored setting restricts the study to pairs in which one side is
always Python, which makes direct dual-execution practical and lets us
isolate the effect of prompt design, direct execution, search, and
evaluator-level aggregation. It does not establish that the same approach
transfers unchanged to arbitrary non-Python language pairs.

A more fundamental limitation is the soundness gap inherent to any
test-based notion of equivalence. Agreement on the tests a trajectory
happens to generate is necessary but not sufficient for genuine semantic
equivalence: two programs that are \emph{not} clones can agree on every
tested input if the model's self-chosen tests do not happen to exercise the
region where they actually differ, producing a false clone; we observed this
concretely in practice, where trajectories that tested only a small number of
similar, narrow, or otherwise low-diversity inputs occasionally concluded
equivalence that a broader or more adversarial set of inputs would have
contradicted. Symmetrically, since the equivalence check between two
executed outputs is made by the model itself rather than by a
language-agnostic, type-aware comparison, superficial cross-language output
differences -- such as numeric formatting or differing integer-overflow
behavior between languages -- can occasionally be read by the model as
evidence of non-equivalence for a pair that is, in fact, a genuine clone. We
therefore treat the framework's verdicts as empirical conclusions over the
specific evidence a given trajectory happened to gather, rather than as a
formal proof of equivalence, and report response rate alongside precision,
recall, and $F_1$ so that abstention is not hidden inside an accuracy number.
